\PassOptionsToPackage{dvipsnames,svgnames,table}{xcolor}
\documentclass[11pt]{article}
\usepackage[a4paper,margin=22mm]{geometry}
\usepackage[no-math]{fontspec}
\setmainfont{DejaVu Serif}
\usepackage{amsmath,amssymb,amsthm,mathtools,graphicx,xcolor,hyperref,xurl,xspace,ifthen,enumerate}
\usepackage[shortlabels]{enumitem}
\setlength{\emergencystretch}{4em}
\AtBeginDocument{\special{pdf:mapline pzdr Dingbats <uzdr.pfb}}
\SetMathAlphabet{\mathbf}{normal}{OT1}{cmr}{bx}{n}
\SetMathAlphabet{\mathbf}{bold}{OT1}{cmr}{bx}{n}
\newtheorem{theo}{Theorem}[section]
\newtheorem{lemm}[theo]{Lemma}
\newtheorem{thm}[theo]{Theorem}
\newtheorem{proposition}[theo]{Proposition}
\newtheorem{corollary}[theo]{Corollary}
\newtheorem{definition}[theo]{Definition}
\newtheorem{theorem}[theo]{Theorem}
\newtheorem*{remas*}{Remarks}
\newtheorem{remas}[theo]{Remarks}
\newtheorem*{exam*}{Example}
\newtheorem{ques}[theo]{Question}
\providecommand{\og}{«\nobreakspace}
\providecommand{\fg}{\nobreakspace»}
\newtheorem{lem}[theo]{Lemma}
\newtheorem{lemma}[theo]{Lemma}
\newtheorem{prop}[theo]{Proposition}
\newtheorem{coro}[theo]{Corollary}
\newtheorem{corol}[theo]{Corollary}
\newtheorem*{conj*}{Conjecture}
\newtheorem{conj}[theo]{Conjecture}
\newtheorem{conjecture}[theo]{Conjecture}
\newtheorem{Proposition}[theo]{Proposition}
\newtheorem{theoreme}[theo]{Theorem}
\theoremstyle{definition}
\newtheorem{defi}[theo]{Definition}
\newtheorem{exam}[theo]{Example}
\newtheorem{exem}[theo]{Example}
\newtheorem{remark}[theo]{Remark}
\newtheorem{example}[theo]{Example}
\newtheorem{exams}[theo]{Examples}
\newtheorem{rema}[theo]{Remark}
\newtheorem{nota}[theo]{Notation}
\newtheorem{notation}[theo]{Notation}
\newtheorem*{notas}{Notations}
\newtheorem{result}[theo]{Result}
\newtheorem{question}[theo]{Question}
\newtheorem{prob}[theo]{Problem}
\newtheorem*{rema*}{Remark}
\newtheorem*{defi*}{Definition}
\newtheorem*{theo*}{Theorem}
\newtheorem*{lemm*}{Lemma}
\newtheorem*{prop*}{Proposition}
\newcommand{\enoncename}{}
\newtheorem{corpusenonce}[theo]{\enoncename}
\newenvironment{enonce}[2][]{\renewcommand{\enoncename}{#2}\begin{corpusenonce}}{\end{corpusenonce}}
\newenvironment{enonce*}[2][]{\par\medskip\noindent\textbf{#2.}\itshape}{\par\medskip}
\newenvironment{proofc}[1][\proofname]{\begin{proof}[#1]}{\end{proof}}
\newcommand{\Mc}{\mathcal{M}}
\newcommand{\pointir}{.\ ---\ }
\newcommand{\equalenv}[2]{\expandafter\let\csname#1\expandafter\endcsname\csname#2\endcsname\expandafter\let\csname end#1\expandafter\endcsname\csname end#2\endcsname}
\providecommand{\diagbox}[3][]{\begin{tikzpicture}[x=1em,y=1em,baseline=(current bounding box.center)]\useasboundingbox(0,0) rectangle(3,1.4);\draw(0,1.4)--(3,0);\node at(.5,.3){#2};\node at(2.5,1.1){#3};\end{tikzpicture}}
\providecommand{\eg}{e.g.\ }
\providecommand{\ie}{i.e.\ }
\providecommand{\cf}{cf.\ }
\providecommand{\longthanks}[1]{\par\smallskip\noindent #1\par\smallskip}
\providecommand{\qbinom}[2]{\genfrac{[}{]}{0pt}{}{#1}{#2}}

\providecommand{\mathof}[1]{\mathrm{#1}}

\numberwithin{equation}{section}


\def\endappendix{}


\usepackage{latexsym,amsmath,amssymb,verbatim,tikz,tikz-cd}
\usepackage{mathtools}
\usepackage{graphicx,xcolor}

\usepackage[english]{babel}
\usepackage{amssymb}
\usepackage{amsmath}
\usepackage{amsthm}
\usepackage{graphicx}
\usepackage{csquotes}
\usepackage{bm}
\usepackage{enumitem}
\usepackage{subcaption}
\usepackage{caption}
\renewcommand{\thesubfigure}{\Alph{subfigure}}
\usepackage{dirtytalk}
\providecommand{\nopunct}{\spacefactor1007\relax}
\newtheorem{problem}[theo]{Problem}

\newcommand{\sym}{\mathrm{Sym}}
\newcommand{\qsym}{\mathrm{QSym}}
\newcommand{\Des}{\operatorname{Des}}
\newcommand{\Asc}{\operatorname{Asc}}
\newcommand{\Syt}{\operatorname{SYT}}
\newcommand{\Short}{\operatorname{Short}}
\newcommand{\Loop}{\operatorname{Loop}}
\newcommand{\Stable}{\operatorname{Stable}}
\newcommand{\len}{\operatorname{len}}
\newcommand{\height}{\operatorname{height}}
\newcommand{\core}{\operatorname{core}}
\newcommand{\res}{\operatorname{res}}
\newcommand{\row}{\operatorname{row}}
\newcommand{\insrt}{\operatorname{insert}}
\newcommand{\UP}{\operatorname{UP}}
\newcommand{\DOWN}{\operatorname{DOWN}}


\newcommand{\ZZ}{\mathbb{Z}}
\newcommand{\NN}{\mathbb{N}}
\newcommand{\RR}{\mathbb{R}}
\newcommand{\QQ}{\mathbb{Q}}
\newcommand{\PP}{\mathbb{P}}

\newcommand{\F}{\mathcal{F}}
\newcommand{\Q}{\mathcal{Q}}
\newcommand{\A}{\mathcal{A}}
\newcommand{\M}{\mathcal{M}}
\newcommand{\I}{\mathcal{I}}
\newcommand{\PM}{\mathcal{PM}}
\newcommand{\Pcal}{\mathcal{P}}


\newcommand{\odds}{\mathrm{odds}}

\makeatletter
\newcommand{\proofstep}[1]{
  \par
  \addvspace{\medskipamount}
  \textbf{#1\@addpunct{.}}\enspace\ignorespaces
}

\newcommand\restr[2]{{
  \left.\kern-\nulldelimiterspace 
  #1 
  \vphantom{\big|} 
  \right|_{#2} 
  }}
\makeatother












\newtheorem{observation}[theo]{Observation}
\title{The reconstruction map is a left inverse of the map given by core and tableau}
\author{Avichai Marmor}
\date{doi:10.5802/alco.351; CC BY 4.0}
\begin{document}
\maketitle
\noindent\textit{Editorial note.} The counted statement is the author's labelled lemma environment at publisher-source.tex lines 982-985; its proof environment opens at line 986, immediately after \textbackslash{}end\{lemma\}, so the pairing is by adjacency and the retained proof range 986-995 is the whole environment.\par\medskip
\noindent\textit{Editorial note.} Two statements are retained as context because the counted proof cites them: lem:properties of core, item (2) (source lines 873-881), used at source line 992, and obs:restriction of matching to set and complement defines uniquely (source lines 692-695), also used at source line 992.\par\medskip
\noindent\textit{Editorial note.} This article carries no label-number dump and the retained counter values are left empty for the whole source, so the printed numbering of the card is the wrapper's default; the identity of the counted item is its author label lem:F\_tilde(F)=id.\par\medskip
\noindent\textit{Editorial note.} r002 note: r001 failed with an undefined reference to def:restriction of matching (addressed by adding that definition as context) and with a LaTeX error because the retained observation context uses the author's observation environment, which the prepared presentation adapter of this source re-declares on the existing wrapper counter; r002 adds that one adapter line as a second preamble span.\par\medskip
\noindent\textit{Editorial scope.} One original result of the pinned publisher article, counted once: the reconstruction map is a left inverse of the map given by core and tableau (author label lem:F\_tilde(F)=id).\par\medskip
\noindent\textit{Allowed standard prerequisites (editorial).}\begin{itemize}
\item The maps F and F-tilde. F sends a matching to the pair formed by its core and its two-row tableau, and F-tilde reconstructs a matching from a core and a tableau.. They are the two maps composed in the statement.
\item Stable set and restriction of a matching. The restriction of a matching to an m-invariant set of vertices is again a matching, and a matching is determined by its restrictions to a set and its complement.. They let the proof reduce the comparison to the stable and the unstable part.
\item Unstable part as a function of the tableau. The unstable chords of a matching are reconstructed from the two-row tableau through the permutation p(T).. It is the equation proved at the start of the argument.
\end{itemize}\medskip
\begin{definition}
\label{def:restriction of matching}
    Given a matching $m \in \M_N$, we say that a set $S \subseteq [N]$ is $m$-invariant if for every chord $(i,j) \in m$ we have $i \in S$ if and only if $j \in S$. Moreover, for a matching $m$ and an $m$-invariant set $S$ we define the \emph{restriction} of $m$ to $S$, denoted $\res_S(m)$, to be a matching on $S$ which is obtained by removing all the vertices not in $S$ from $m$.
\end{definition}

\begin{lemma}
\label{lem:properties of core}
    The reduction process has the following properties:
    \begin{enumerate}[label=(\arabic*), ref=\ref{lem:properties of core}\text{ (part~}\arabic*)]
        \item \label{lem:properties of core--intersect stable} If a chord intersects another chord then it is stable.
        \item \label{lem:properties of core--chord unstable iff all midpoints unstable} A chord $(i,j)$ is stable if and only if there exists a stable vertex in $\{i+1,\dots,j-1\}$.
        \item \label{lem:properties of core--perfect non-crossing unstable} Given a matching $m$ and $i < j$, if the set $[i,j]$ is $m$-invariant \textup{(}as defined in Definition~\ref{def:restriction of matching}\textup{)} and the restricted matching $\res_{[i,j]}(m)$ is perfect and non-crossing, then $\ell \notin \Stable(m)$ for all $\ell \in [i,j]$.
    \end{enumerate}
\end{lemma}

\begin{observation}
\label{obs:restriction of matching to set and complement defines uniquely}
    Let $S \subseteq [N]$ be a set of vertices, let $m_1$ be a matching on $S$, and let $m_2$ be a matching on $[N]\setminus S$. Then there exists a unique matching $m \in \M_N$ such that $\res_S(m) = m_1$ and $\res_{[N] \setminus S}(m) = m_2$.
\end{observation}

\begin{lemma}
\label{lem:F_tilde(F)=id}
    Let $m \in \M_{N,f}$ be a matching, and denote $F(m) = (\core(m),T)$. Then $\tilde{F}(\core(m),T) = m$.
\end{lemma}

\begin{proof}
    Denote $\left|\Stable(m)\right| = N-2k$, implying that $\core(m) \in \M_{N-2k,f}$ and $T \in \Syt(N-k,k)$. Moreover, denote $p=p(T) \in \Pcal_{N,N-2k}$. First, we prove that
    \begin{equation}
    \label{eqn:lemma F_tilde(F)=id}
        \res_{[N]\setminus\Stable(m)}(m) = m_{\text{unstable}}(p).
    \end{equation}
    We note that the set $[N] \setminus \Stable(m)$ is $m$-invariant, so the left-hand side of Equation~\eqref{eqn:lemma F_tilde(F)=id} is well-defined. Notice that $|[N]\setminus\Stable(m)| = |[N]\setminus\Stable(p)| = 2k$ and that $\res_{[N]\setminus\Stable(m)}(m)$ is a perfect matching. Therefore, in order to prove that Equation~\eqref{eqn:lemma F_tilde(F)=id} holds, it suffices to prove that any unstable chord of $m$ belongs to $m_{\text{unstable}}(p)$ too. Let $(i,j) \in m$ be an unstable chord. Thus, we may deduce that $i \in \UP(p)$ and $j \in \DOWN(p)$. Let $i' \in \UP(p)$ such that $i < i' < j$. By Lemma~\ref{lem:properties of core--chord unstable iff all midpoints unstable}, $i'$ is an endpoint of an unstable chord $(i',j') \in m$ with $i' < j'$. By Lemma~\ref{lem:properties of core--intersect stable}, we obtain that the chord $(i,j)$ does not intersect $(i',j')$, implying that $i < i' < j' < j$. On the other hand, every $j' \in \DOWN(p)$ with $i < j' < j$ is an endpoint of an unstable chord $(i',j') \in m$ with $i < i' < j' < j$. Therefore, we obtain a bijection from $\UP(p) \cap [i,j]$ to $\DOWN(p) \cap [i,j]$ (with $i$ matched with $j$), so $\height_p(i) = \height_p(j)$. Moreover, this bijection has the property that if $i' \in \UP(p) \cap [i,j]$ is matched with $j' \in \DOWN(p) \cap [i,j]$ then $i' < j'$ (\ie the ascending steps of the path appear before the associated descending steps), so $\height_p(j) < \height_p(i')$ for all $i' \in \UP(p) \cap [i+1,j-1]$. We may conclude that $(i,j) \in m_{\text{unstable}}(p)$ for every unstable chord $(i,j)$ of $m$ and prove that Equation~\eqref{eqn:lemma F_tilde(F)=id} holds.

    Next, we denote $m' = \tilde{F}(\core(m),T)$ and prove that $m=m'$. From Equation~\eqref{eqn:lemma F_tilde(F)=id} we deduce that the supports of the matchings $\res_{[N]\setminus\Stable(m)}(m)$ and $m_{\text{unstable}}(p)$ are identical, and therefore $\Stable(m) = \Stable(p)$. Thus, the set $[N]\setminus\Stable(m)$ is both $m$-invariant and $m'$-invariant, and restricting each of these matchings to $[N]\setminus\Stable(m)$ results in $m_{\text{unstable}}(p)$. In addition, it can be easily verified from the descriptions of $F$ and $\tilde{F}$ that $\res_{\Stable(m)}(m) = \res_{\Stable(m)}(m')$ is the matching obtained by relabeling the vertices of $\core(m)$ with the elements of $\Stable(m)$ in increasing order. By Observation~\ref{obs:restriction of matching to set and complement defines uniquely}, we may deduce that $m=m'$.
\end{proof}

\end{document}
